\documentclass[11pt]{article}
\usepackage[T1]{fontenc}
\usepackage{lmodern,amsmath,amssymb,amsthm,booktabs,tabularx}
\usepackage[margin=1in]{geometry}
\usepackage{microtype}
\usepackage[colorlinks=true,linkcolor=blue,citecolor=blue,urlcolor=blue]{hyperref}
\newcommand{\Tr}{\operatorname{Tr}}
\newcommand{\Ree}{\operatorname{Re}}
\newcommand{\rank}{\operatorname{rank}}
\newcommand{\C}{\mathbb C}
\newcommand{\N}{\mathbb N}
\newtheorem{theorem}{Theorem}
\newtheorem{lemma}[theorem]{Lemma}
\newtheorem{corollary}[theorem]{Corollary}
\title{Sharp accretive modulus bounds from Hayashi's (2001) pinching inequality:\\
extending Abed--Nikolaeva--Novikov (2026),\\
arXiv:2610.34388v1, Section 11.2(5) and Remark 5.45}
\author{Haobo Ma\\{\small ChronoAI Pte Ltd, Singapore}
\and Wenlin Zhang\\{\small National University of Singapore}\\
{\small ChronoAI Pte Ltd, Singapore}}
\date{6 November 2026}
\begin{document}
\maketitle
\begin{abstract}
Applying Hayashi's (2001) pinching inequality and its projection-count form
in Ogawa--Hayashi (2002), we extend Abed--Nikolaeva--Novikov (2026),
arXiv:2610.34388v1, Proposition 5.44(c), and answer the density-independent
boundedness question in Section 11.2(5).
For every accretive matrix $X$, we obtain $\Ree X\leq m(X)|X|$,
where $m(X)$ counts its distinct positive singular values.
For a positive definite density of condition number $\kappa$, this gives
$1\leq a(\rho)\leq n\kappa/(\kappa+n-1)<n$ when $n>1$.
The optimal operator constant in Remark 5.45 is $c_n=n$, and
$\sup_{\rho>0}a(\rho)=n$, also with rank-one Hermitian parts.
The sharpness mechanism is classical: the Liv\v{s}ic triangular model
displayed by Arlinskii--Tsekanovskii (2006) supplies such examples;
we also give rational strictly accretive witnesses.
\end{abstract}

\section{The question and the matrix bound}

For $X\in M_n(\C)$ put $|X|=(X^*X)^{1/2}$ and
$\Ree X=(X+X^*)/2$. An accretive matrix has $\Ree X\geq0$.
For $\rho>0$, write $\varphi_\rho(Y)=\Tr(\rho Y)$ and define
\begin{equation}\label{eq:def}
a(\rho)=\sup_{\substack{A\geq0,\ B=B^*\\A+iB\ne0}}
\frac{\Tr(\rho A)}{\Tr(\rho|A+iB|)}.
\end{equation}
These are the conventions of \cite[Sections 3.1 and 5.12]{ANN};
no trace normalization is imposed on $\rho$.
Section 11.2(5), printed page 59, poses the question
``Is $a$ bounded on $M_n(\C)$ independently of $\rho$?''
Here $n$ is fixed. Proposition 5.44(c) of \cite{ANN} already gives
$\Ree X\leq2|X|$ in dimension two and $a(\rho)<2$ there.
The argument below gives the all-dimensional answer, a singular-value
refinement, and the optimal dimension constant of Remark 5.45.
The underlying positive-matrix estimate is Hayashi's pinching inequality
\cite[Lemma 5]{H}; its projection-count version appears in
\cite[Appendix B, Lemma 4, (61)]{OH}.
The similarity-to-unitary and trace method also appears in fedja's
2017 answer \cite{fedja}, in a different operator comparison.
These sources provide the pinching and similarity methods; their
accretive-modulus specialization is the argument below. Winter (2025)
\cite[Section I, Lemma 1, (4)]{W} gives weighted generalisations of the
same pinching inequality.

The earlier bound in \cite[Proposition 5.44(a)]{ANN} is
$a(\rho)\leq G(\kappa)=(\kappa+1)/(2\sqrt\kappa)$.
It remains applicable alongside the dimension-dependent estimate here:
their minimum is an upper bound. The uniform ceiling $n$, rather than
a pointwise improvement of $G(\kappa)$, answers the fixed-dimensional
boundedness question.

For comparison, the modulus estimates of Furuichi, Moradi and Sababheh
\cite[Section 2.5, Theorem 2.6]{FMS} assume a numerical range in a fixed
sector of half-angle $\alpha<\pi/2$ and contain a factor $\sec\alpha$.
Here the comparison bounds $\Ree X$ by $|X|$ with a coefficient depending
only on dimension and includes a singular Hermitian part.

\paragraph{Literature scope.}
Full-text access to Moradi--Furuichi--Sababheh (2026) \cite{MFS26}
was unavailable; comparison is limited to its abstract and bibliography.
Whether its theorems or factorisations establish or directly imply the
dimension or distinct-singular-value Loewner-order bounds, the weighted
density estimate, or the sharp coefficient $c_n=n$ remains unresolved.
The abstract's sector-angle terminology alone does not exclude those
overlaps. The results below are explicit deductions from the credited
classical methods, without a priority assertion relative to that article.

\begin{lemma}\label{lem:matrix}
If $X\in M_n(\C)$ is accretive, let $m(X)$ be the number of distinct
positive singular values, with $m(0)=0$. Then
\begin{equation}\label{eq:matrix}
0\leq\Ree X\leq m(X)|X|\leq\rank(X)|X|\leq n|X|.
\end{equation}
\end{lemma}
\begin{proof}
Set $A=\Ree X$, $P=|X|$, and $Q=A^{1/2}$.
Write $P=\sum_{j=1}^m s_jE_j$ for its distinct positive eigenvalues
and their orthogonal spectral projections; let $E_0$ project onto $\ker P$.
For $v$ in the range of $E_j$, $\|Xv\|=s_j\|v\|$, so
\begin{equation}\label{eq:diag}
0\leq\langle v,Av\rangle=\Ree\langle v,Xv\rangle
\leq s_j\|v\|^2.
\end{equation}
The same argument on $\ker P=\ker X$ gives $QE_0=0$.
The triangle and Cauchy--Schwarz inequalities therefore give, for any $v$,
\begin{align*}
\langle v,Av\rangle
&=\left\|\sum_{j=1}^m QE_jv\right\|^2
\leq m\sum_{j=1}^m\|QE_jv\|^2\\
&\leq m\sum_{j=1}^m s_j\|E_jv\|^2
=m\langle v,Pv\rangle.
\end{align*}
For $X=0$ the assertion is immediate. This is the projection-count
pinching proof, including zero singular values.
\end{proof}

\begin{corollary}\label{cor:condition}
Let $\rho>0$ have extreme eigenvalues $L\geq l>0$ and condition
number $\kappa=L/l$. For accretive $X\ne0$, with $m=m(X)$,
\[
\frac{\Tr(\rho\Ree X)}{\Tr(\rho|X|)}
\leq\frac{m\kappa}{\kappa+m-1}.
\]
\end{corollary}
\begin{proof}
Put $A=\Ree X$, $P=|X|$, $N=\Tr(\rho A)$ and $D=\Tr(\rho P)>0$.
Summing the diagonal estimate \eqref{eq:diag} in a singular-vector
basis gives $\Tr A\leq\Tr P$.
The positive matrix $mP-A$ thus has trace at least $(m-1)\Tr P$.
Consequently
\[
mD-N\geq l(m-1)\Tr P\geq \frac{l(m-1)}{L}N,
\]
where $N\leq L\Tr A\leq L\Tr P$. Rearranging proves the claim.
\end{proof}

\begin{theorem}\label{thm:density}
For every $n\geq1$ and $\rho\in M_n(\C)$ with $\rho>0$,
put $\kappa=\lambda_{\max}(\rho)/\lambda_{\min}(\rho)$.
The set in \eqref{eq:def} is nonempty and has a real least upper bound,
and
\[
1\leq a(\rho)\leq\frac{n\kappa}{\kappa+n-1}\leq n.
\]
If $n=1$, then $a(\rho)=1$. If $n>1$, then $a(\rho)<n$.
The supremum is attained for every $\rho>0$.
\end{theorem}
\begin{proof}
For any nonzero $P\geq0$, if $m=\lambda_{\min}(\rho)>0$, then
\[
\Tr(\rho P)\geq m\Tr P>0.
\]
Here the nonnegative trace pairing follows from
$\Tr(RP)=\Tr(P^{1/2}RP^{1/2})\geq0$ for $R,P\geq0$.
Consequently every denominator in \eqref{eq:def} is positive.
Taking trace pairings in \eqref{eq:matrix} bounds every quotient
between zero and $n$. The choice $A=I$, $B=0$ gives the value one.
Completeness of the real numbers supplies the least upper bound.
Corollary~\ref{cor:condition} and monotonicity of
$u\kappa/(\kappa+u-1)$ for $u\geq1$ give the stated bound, including
strictness for $n>1$. For $n=1$ the bounds coincide.

For attainment, the compactness argument of
\cite[Proposition 5.44(b)]{ANN} applies: restrict by homogeneity to
\[
\{X:\Ree X\geq0,\ \Tr(X^*X)=1\}.
\]
This is a nonempty compact set. On it,
$\Tr(\rho|X|)\geq m\Tr|X|\geq m$, since the sum of singular values
is at least their Euclidean norm. The quotient is continuous, because
the positive square root is continuous in finite dimension, and so
attains its maximum.
\end{proof}

\begin{corollary}\label{cor:positive}
Restricting \eqref{eq:def} to $B\geq0$ gives a constant $a_+(\rho)$
with $1\leq a_+(\rho)\leq a(\rho)\leq n\kappa/(\kappa+n-1)$.
For $n>1$ the last expression is strictly less than $n$.
The restricted supremum is attained.
\end{corollary}
\begin{proof}
The restricted set is a subset of the original set and contains
the choice $A=I$, $B=0$. Its Frobenius-normalized set is closed
in the compact set above; the same continuity argument applies.
\end{proof}

\section{Sharpness in every dimension}

\begin{theorem}\label{thm:sharp}
The least constant $c_n$ such that $\Ree X\leq c_n|X|$ for all
accretive $X\in M_n(\C)$ is $n$. More generally the optimal coefficient
for matrices of fixed rank $1\leq r\leq n$ is $r$. Moreover,
\[
\sup_{\rho\in M_n(\C),\ \rho>0}a(\rho)=n.
\]
The same supremum holds if $\Tr\rho=1$ is required.
\end{theorem}
\begin{proof}
The upper bounds are already proved. The following strictly accretive
triangular family implements the similarity/congruence mechanism in
\cite[following Proposition 5.44]{ANN}; the classical rank-one
Hermitian-part version is given below. Fix $0<t<1/n$, set
$h_{n+1}=0$, and for $1\leq j\leq n$ put
\[
\lambda_j=\frac{1+ijt}{1-ijt},\qquad
h_j=\Ree\lambda_j=\frac{1-j^2t^2}{1+j^2t^2}.
\]
The numbers $\lambda_j$ are distinct and have modulus one, and
$h_1>\cdots>h_n>0$. Define an upper triangular matrix $T$ by
\[
T_{jj}=\lambda_j,\qquad T_{ij}=2h_j\ (i<j),\qquad T_{ij}=0\ (i>j).
\]
Its Hermitian part $H=\Ree T$ satisfies $H_{ij}=h_{\max(i,j)}$.
If $b_k$ is the vector with its first $k$ entries equal to one and
the rest zero, then
\begin{equation}\label{eq:gram}
H=\sum_{k=1}^n(h_k-h_{k+1})b_kb_k^*>0.
\end{equation}
The positivity is strict because all coefficients are positive and
$b_1,\ldots,b_n$ form a basis.

The distinct eigenvalues make $T$ diagonalizable: write $T=SDS^{-1}$
with $D=\operatorname{diag}(\lambda_1,\ldots,\lambda_n)$.
Set $G=(S^{-1})^*S^{-1}>0$. Since $D^*D=I$, one has $T^*GT=G$.
For
\[
X=G^{1/2}TG^{1/2},
\]
it follows that
$X^*X=G^{1/2}T^*GTG^{1/2}=G^2$,
so $|X|=G$ and $\Ree X=G^{1/2}HG^{1/2}>0$.
Let $w=(1,\ldots,1)^\mathsf{T}$, choose $\delta>0$, and set
\[
W=ww^*+\delta I>0,\qquad
\rho=G^{-1/2}WG^{-1/2}>0.
\]
Cyclicity of the trace gives the admissible quotient
\begin{equation}\label{eq:sharpquotient}
\frac{\Tr(\rho\Ree X)}{\Tr(\rho|X|)}
=\frac{\Tr(WH)}{\Tr W}
=\frac{\displaystyle\sum_{i,j=1}^n h_{\max(i,j)}
      +\delta\sum_{j=1}^n h_j}{n(1+\delta)}.
\end{equation}
As $t\downarrow0$, $H$ tends to $ww^*$, so this expression tends to
$(n+\delta)/(1+\delta)$. Letting $\delta\downarrow0$ gives $n$.
In particular no smaller universal matrix constant can hold, since
such a constant would bound all the trace quotients above.
The rank-$r$ statement follows by taking this construction in dimension
$r$ and adjoining a zero block. Finally $a(s\rho)=a(\rho)$ for every $s>0$, by cancellation in
\eqref{eq:def}, so normalizing $\rho$ does not change the supremum.
\end{proof}

\begin{corollary}\label{cor:rankone}
Both sharp suprema in Theorem~\ref{thm:sharp} remain $n$ under the
restriction $\rank(\Ree X)=1$.
\end{corollary}
\begin{proof}
In the notation above replace the strict upper entries of $T$ by
$2\sqrt{h_ih_j}$ for $i<j$. Then $\Ree T=vv^*$, where
$v=(\sqrt{h_1},\ldots,\sqrt{h_n})^{\mathsf T}$, and
$\|\Ree T\|=\sum_jh_j\to n$ as $t\downarrow0$.
This is the Liv\v{s}ic triangular model in
\cite[Section 2.3, Theorem 2.7]{AT}, multiplied by $-i$:
take $\alpha_j=-\Im\lambda_j$, $\beta_j=\sqrt{2h_j}$ there.
Its distinct unimodular eigenvalues give $G$ and $X$ as above.
Congruence preserves the rank of the Hermitian part.
For $q=v/\|v\|$ and $W=qq^*+\delta I$, the corresponding density
$\rho=G^{-1/2}WG^{-1/2}>0$ has quotient
\[
\frac{\Tr(Wvv^*)}{\Tr W}
=\left(\sum_jh_j\right)\frac{1+\delta}{1+n\delta}\longrightarrow n
\]
as $t,\delta\downarrow0$.
\end{proof}

\begin{corollary}\label{cor:strictoperator}
For each accretive matrix of rank $r>1$ there is a coefficient
$c(X)<r$ with $\Ree X\leq c(X)|X|$. The optimal rank coefficient
$r$ is therefore a supremum, not an attained operator ratio.
\end{corollary}
\begin{proof}
If $Xv=0$ then $\langle v,\Ree Xv\rangle=0$, so
$(\Ree X)v=0$ and $X^*v=0$. Thus $\ker X$ reduces $X$.
Restrict to its orthogonal complement, where $P=|X|>0$ and
$X=UP$ with $U$ unitary. The matrix
$T=P^{-1/2}XP^{-1/2}=P^{-1/2}UP^{1/2}$ is similar to $U$.
Consequently $H=\Ree T\geq0$ has
$\lambda_{\max}(H)\leq\Tr H=\Ree\Tr U\leq r$.
Equality $\lambda_{\max}(H)=r$ would force every eigenvalue of $U$
to be $1$, hence $U=I$, $T=I$ and $H=I$, impossible for $r>1$.
Take $c(X)=\lambda_{\max}(H)$ and extend by zero on $\ker X$.
This is the similarity-and-trace method of \cite{fedja};
the dimension-two nonattainment appears in \cite{ANN}.
\end{proof}

\begin{corollary}\label{cor:blocks}
For positive definite $\rho_1,\rho_2$,
$a(\rho_1\oplus\rho_2)\geq\max\{a(\rho_1),a(\rho_2)\}$.
The same statement holds for $a_+$.
\end{corollary}
\begin{proof}
Use $X=X_1\oplus0$ and $X=0\oplus X_2$ respectively in
\eqref{eq:def}; the quotients are exactly the smaller-block quotients.
The condition on the imaginary part is preserved.
\end{proof}

\section{Coverage and verification}

All locators in this table refer to arXiv:2610.34388v1 \cite{ANN}.
Questions are labelled ``Answered'', assertions already proved in the
source ``Reproved'', and additional conclusions ``Not posed; proved''.
\begin{center}
\small
\begin{tabularx}{\textwidth}{@{}>{\raggedright\arraybackslash}p{0.32\textwidth}>{\raggedright\arraybackslash}X@{}}
\toprule
Source item & Coverage\\
\midrule
Section 11.2(5), boundedness on $M_n(\C)$ &
Answered: yes, with $a(\rho)\leq n\kappa/(\kappa+n-1)$
(Theorem~\ref{thm:density}).\\
Section 11.2(5), same boundedness with $B\geq0$ &
Answered: yes (Corollary~\ref{cor:positive}); the same strengthened upper bound.\\
Remark 5.45, best operator constant $c_n$ &
Not posed as an equality; proved: $c_n=n$
(Theorem~\ref{thm:sharp}), settling the numerical lower bound $c_3\geq2.94$.
Reproved; prior proof of $c_2=2$: Abed--Nikolaeva--Novikov (2026),
Proposition 5.44(c) and its following paragraph.\\
Proposition 5.44(b),(c), attainment and the $n=2$ upper bound &
Reproved; prior proof Abed--Nikolaeva--Novikov (2026).
Theorem~\ref{thm:density} also gives the all-dimensional strict bound.\\
Section 11.2(5), determine $a(\rho)$ and $a_+(\rho)$ &
Not covered: the full density-dependent formulas; Theorems~\ref{thm:density}
and \ref{thm:sharp} give bounds and the supremum over densities.\\
Section 11.2(5), II$_1$-factor equality and its $B\geq0$ variant &
Not covered.\\
Remark 5.45, rank-one maximisers &
Not covered for a fixed density. Not posed; proved: rank-one Hermitian
parts suffice for the global sharp coefficient (Corollary~\ref{cor:rankone}).\\
Remark 5.45, extreme-density multiplicities &
Not covered: exact dependence or independence. Not posed; proved:
block monotonicity (Corollary~\ref{cor:blocks}).\\
Remark 5.45, numerical values, $\kappa\to1$ ratio and positive-$B$ values &
Not covered; no numerical equality or asymptotic expansion is inferred.\\
Singular-value, condition-number and rank refinements &
Not posed; proved: Lemma~\ref{lem:matrix}, Corollary~\ref{cor:condition},
Theorem~\ref{thm:sharp} and Corollary~\ref{cor:strictoperator}.
The pinching theorem is prior work of Hayashi (2001), also
Ogawa--Hayashi (2002); the accretive specialization is supplied above.\\
\bottomrule
\end{tabularx}
\end{center}

The declarations in the
\href{https://github.com/the-omega-institute/trureturing/blob/d9bb8a406783c5222849b798f3c2ab738f653d75/D5/S3/Quantum/Matrix/AbedNikolaevaNovikovAFDensityBound.lean}{pinned Lean source}
have namespace
\begin{quote}\small\ttfamily
D5.S3.Quantum.Matrix.\\
AbedNikolaevaNovikovAFDensityBound
\end{quote}
The declaration \texttt{accretive\_le\_card\_abs} proves
$A\leq n|A+iB|$ for $A\geq0$, $B=B^*$, including singular $A+iB$.
The imported
\href{https://github.com/the-omega-institute/trureturing/blob/d9bb8a406783c5222849b798f3c2ab738f653d75/D5/S3/Quantum/FiniteDimensional.lean}{finite-dimensional trace definition}
is \texttt{bornProbability rho P = Matrix.trace (rho * P)};
it imposes no trace-one normalization.
The declaration \texttt{result} proves that the exact quotient set is
nonempty, has its stated least upper bound, and admits the
density-independent witness $C=n$ for every $n\geq1$.
The distinct-singular-value and rank refinements, condition-number bound,
strictness, attainment, positive-$B$ deduction, block monotonicity and
sharpness have written proofs above; they are not conclusions of these
two declarations.
The accompanying \texttt{verify.py} uses exact rational arithmetic to
check noncommuting and singular matrix examples in dimensions $1$--$6$
and the rational sharpness construction in dimensions $1$--$8$,
including its diagonalizer and the identity $T^*GT=G$ over $\mathbb Q(i)$.
These finite checks illustrate the identities; the proofs establish
the statements for all dimensions.

\begin{thebibliography}{8}
\bibitem{ANN}
Sami Abdullah Abed, Irina Nikolaeva and Andrej Novikov,
\emph{How characterisations of traces and central elements fail:
sharp violation constants via Michelson contrast and noncentrality},
arXiv:2610.34388v1 (25 October 2026), Sections 3.1, 5.12 and 11.2(5).
\href{https://arxiv.org/abs/2610.34388v1}{arXiv:2610.34388v1}.
\bibitem{H}
Masahito Hayashi,
\emph{Optimal sequence of POVMs in the sense of Stein's lemma in
quantum hypothesis testing}, arXiv:quant-ph/0107004v2 (9 August 2001),
Section 3, Lemma 5, p.~6.
\href{https://arxiv.org/abs/quant-ph/0107004v2}{quant-ph/0107004v2}.
\bibitem{OH}
Tomohiro Ogawa and Masahito Hayashi,
\emph{On error exponents in quantum hypothesis testing},
arXiv:quant-ph/0206151v3 (14 August 2002), Appendix B,
Lemma 4, (61), pp.~10--11.
\href{https://arxiv.org/abs/quant-ph/0206151v3}{quant-ph/0206151v3}.
\bibitem{AT}
Yury Arlinskii and Eduard Tsekanovskii,
\emph{Non-self-adjoint Jacobi matrices with rank one imaginary part},
arXiv:math/0602033v1 (2 February 2006), Section 2.3, Theorem 2.7, p.~10.
\href{https://arxiv.org/abs/math/0602033v1}{math/0602033v1}.
\bibitem{fedja}
fedja, answer to \emph{Comparison of the absolute value of an operator
with its positive parts}, MathOverflow, 8 November 2017, answer 285512,
first three paragraphs.
\href{https://mathoverflow.net/a/285512}{MathOverflow answer 285512}.
\bibitem{FMS}
Shigeru Furuichi, Hamid Reza Moradi and Mohammad Sababheh,
\emph{Further properties of accretive matrices},
Annales Fennici Mathematici \textbf{49} (2024), 387--404.
\href{https://doi.org/10.54330/afm.146278}{doi:10.54330/afm.146278}.
\bibitem{W}
Andreas Winter,
\emph{Generalised Pinching Inequality},
arXiv:2510.19111v1 (21 October 2025), Section I, Lemma 1, (4).
\href{https://arxiv.org/abs/2510.19111v1}{arXiv:2510.19111v1}.
\bibitem{MFS26}
Hamid Reza Moradi, Shigeru Furuichi and Mohammad Sababheh,
\emph{Norm bounds and factorization results for accretive matrices},
The Journal of Analysis \textbf{34} (2026), 2641--2660,
published online 8 June 2026.
\href{https://doi.org/10.1007/s41478-026-01123-8}{doi:10.1007/s41478-026-01123-8}.
\end{thebibliography}
\medskip
\noindent\small\textcopyright\ 2026 Haobo Ma and Wenlin Zhang.
Licensed under \href{https://creativecommons.org/licenses/by/4.0/}{CC BY 4.0}.
\end{document}
